\section{Integration with the Hanzo Agent Harness}

\label{sec:harness}

\subsection{The Continuous Learning Extension}

FAI is implemented as a \emph{plugin} for the Hanzo agent harness,
augmenting the base agent with four capabilities:
\begin{enumerate}[leftmargin=1.5em]
    \item \textbf{Session telemetry collection}: Intercept and log all tool
          calls and outcomes without modifying agent behavior.
    \item \textbf{Background GRPO}: Run experience extraction asynchronously
          between sessions, not blocking agent execution.
    \item \textbf{DSO network tools}: Expose DSO operations as callable tools
          so the agent can query its own fleet status.
    \item \textbf{Periodic synchronization}: Background service that runs
          library consolidation and DSO broadcast on a configurable schedule.
\end{enumerate}

\subsection{DSO Network Tools}

The plugin exposes four tools to the agent:

\begin{verbatim}
dso_network(action="status")
  -> { peers: N, library_size: M, last_sync: timestamp,
       bytes_sent: X, bytes_received: Y }

dso_network(action="broadcast", domain="coding")
  -> { experiences_sent: K, peers_reached: P, batch_id: str }

dso_network(action="receive", domain="coding", limit=50)
  -> { experiences_received: K, library_delta: D }

dso_network(action="peers")
  -> { peers: [{ node_id, latency_ms, reputation, last_seen }] }
\end{verbatim}

These tools allow the agent to actively manage its DSO participation---for
example, by broadcasting immediately after completing a difficult task, or
by requesting experiences from a specific domain before starting a new
project.

\subsection{Experience Lifecycle}

The complete lifecycle of an experience from agent session to fleet benefit
is illustrated in Figure~\ref{fig:lifecycle} and described below:

\begin{enumerate}[leftmargin=1.5em]
    \item \textbf{Session}: Agent executes tools, producing telemetry
          $\{(\text{tool}_t, \text{input}_t, \text{output}_t, r_t)\}$.
    \item \textbf{Trajectory construction}: Telemetry is grouped into
          trajectories $\tau = (\tau_1, \ldots, \tau_T)$ by task.
    \item \textbf{GRPO extraction}: Group relative advantages are computed
          and semantic patterns are extracted via LLM introspection.
    \item \textbf{Local library update}: Patterns are embedded and stored
          in $\mathcal{L}_i$, indexed by task embedding.
    \item \textbf{BitDelta compression}: Library is compressed using
          Algorithm~\ref{alg:compress} with the current baseline.
    \item \textbf{DSO broadcast}: Compressed batch is broadcast to $k$
          peers via the gossip protocol.
    \item \textbf{Peer receive}: Peers receive the batch, verify integrity,
          and add to their incoming queue.
    \item \textbf{DeltaSoup merge}: Algorithm~\ref{alg:deltasoup} merges
          all received batches into the peer's local library.
    \item \textbf{Better next session}: The updated library is used for
          PoE decoding (Eq.~\ref{eq:poe}) in the peer's next session.
\end{enumerate}

\begin{figure}[H]
\centering
\small
\setlength{\tabcolsep}{4pt}
\begin{tabular}{ccccccccc}
\textbf{Session} & $\to$ & \textbf{Telemetry} & $\to$ & \textbf{GRPO} & $\to$ & \textbf{Local} & $\to$ & \textbf{Compress} \\
& & & & & & $\mathcal{L}_i$ & & BitDelta \\[4pt]
\multicolumn{3}{c}{\textbf{Better Session}} & $\leftarrow$ & \textbf{DeltaSoup} & $\leftarrow$ & \textbf{Receive} & $\leftarrow$ & \textbf{Broadcast} \\
\multicolumn{3}{c}{(PoE decoding)} & & merge & & peers & & DSO gossip \\
\end{tabular}
\caption{Experience lifecycle: from agent session to fleet benefit and back.}
\label{fig:lifecycle}
\end{figure}

\subsection{Interaction with Self-Improvement}

The Hanzo agent harness includes a self-improvement loop that periodically
audits agent behavior and proposes policy changes. FAI integrates with this
loop at two points:

\begin{enumerate}[leftmargin=1.5em]
    \item \textbf{Library statistics inform sharing policy}: The self-improvement
          loop's maintenance phase (Loop 4) examines experience library statistics
          (domain coverage, quality distribution, staleness) and can increase
          broadcast frequency for under-represented domains or reduce it for
          saturated ones.
    \item \textbf{Fleet quality informs individual improvement}: Low fleet-wide
          quality in a domain signals that new local GRPO runs may be productive.
          The self-improvement loop uses $Q^{(T)}$ as a signal for when to
          invest in fresh rollouts.
\end{enumerate}

% ============================================================================
% 8. PRIVACY AND SECURITY
% ============================================================================
